% ============================================================
% Title supplied by appendix wrapper.
\label{sec:advanced-operators-iii}
% ============================================================

These tools extend the operator toolbox beyond the main spine's transport and
monodromy interfaces by adding principal angles (subspace mismatch via SVD),
invariant filters (projection onto admissible subalgebras), and unitary
splittings (norm bookkeeping across observable and shadow components).

% ------------------------------------------------------------
\subsection{Principal Angles Between Invariant Subspaces (SVD Machinery)}
\label{subsec:principal-angles}
% ------------------------------------------------------------

Principal angles measure the ``mismatch'' or ``frustration'' between two subspaces, encoding how far they are from being orthogonal or coincident.

\begin{definition}[Principal angles]
\label{def:principal-angles}\lean{CoherenceCalculus.AppendixOperators.principalAngles}
Let $U$ and $V$ be subspaces of a finite-dimensional inner-product space $\cH$, with $\dim U = p$ and $\dim V = q$ where $p \le q$.
The \emph{principal angles} $0 \le \theta_1 \le \theta_2 \le \cdots \le \theta_p \le \pi/2$ are defined recursively:
\[
\cos\theta_k := \max_{\substack{u \in U,\, v \in V \\ \|u\|=\|v\|=1 \\ u \perp u_1, \ldots, u_{k-1} \\ v \perp v_1, \ldots, v_{k-1}}} \langle u, v \rangle,
\]
where $(u_k, v_k)$ is a pair achieving the maximum at step $k$.
\end{definition}

\begin{definition}[Principal angle operator]
\label{def:principal-angle-operator}\lean{CoherenceCalculus.AppendixOperators.principalAngleOp}
Let $P_U$ and $P_V$ be the orthogonal projections onto $U$ and $V$ respectively.
Define the \emph{principal angle operator}
\[
A_{UV} := P_U P_V|_V : V \to U,
\]
the restriction of $P_U P_V$ to $V$.
\end{definition}

\begin{proposition}[SVD characterization of principal angles]
\label{prop:svd-principal}\lean{CoherenceCalculus.AppendixOperators.svd_principal_angles}
The singular values of $A_{UV} = P_U P_V|_V$ are $\sigma_k = \cos\theta_k$, where $\theta_k$ are the principal angles between $U$ and $V$.
\end{proposition}

\begin{proof}
Choose orthonormal bases of \(U\) and \(V\). In those bases, the operator
\(A_{UV} = P_U P_V|_V\) is represented by the matrix
\[
M_{ij} = \langle u_i, v_j \rangle .
\]
Apply the singular value decomposition to \(M\). There exist orthonormal bases
\(\{\tilde u_i\}\) of \(U\) and \(\{\tilde v_j\}\) of \(V\) such that
\[
\langle \tilde u_i,\tilde v_j\rangle = \sigma_i \delta_{ij}
\]
for singular values \(\sigma_1 \ge \cdots \ge \sigma_p \ge 0\).

The recursive maximization defining the principal angles selects exactly these
orthonormal pairs, and the maximizing value at step \(i\) is
\(\cos\theta_i = \sigma_i\). Since the singular values of \(A_{UV}\) are the
singular values of \(M\), the claim follows.
\end{proof}

\begin{definition}[Frustration tensor]
\label{def:frustration-tensor}\lean{CoherenceCalculus.AppendixOperators.frustrationTensor}
The \emph{frustration tensor} (or \emph{overlap tensor}) between subspaces $U$ and $V$ is the SVD data of $A_{UV}$:
\[
\mathcal{F}(U, V) := \bigl( \sigma_1, \sigma_2, \ldots, \sigma_{\min(p,q)} \bigr) = \bigl( \cos\theta_1, \cos\theta_2, \ldots, \cos\theta_{\min(p,q)} \bigr).
\]
Equivalently, $\mathcal{F}(U,V)$ can be represented as the matrix of inner products $\langle u_i, v_j \rangle$ in principal bases.
\end{definition}

\begin{corollary}[Extremal cases]
\label{cor:principal-extremes}\lean{CoherenceCalculus.AppendixOperators.subspace_inclusion_iff_angles_zero}\lean{CoherenceCalculus.AppendixOperators.subspace_orthogonal_iff_angles_max}\lean{CoherenceCalculus.AppendixOperators.subspace_equality_iff_all_ones}
\begin{enumerate}[label=(\roman*)]
\item $U \subseteq V$ if and only if all principal angles are zero (all $\sigma_k = 1$).
\item $U \perp V$ if and only if all principal angles are $\pi/2$ (all $\sigma_k = 0$).
\item $U = V$ (when $\dim U = \dim V$) if and only if all $\sigma_k = 1$.
\end{enumerate}
\end{corollary}

\begin{proof}
Direct from the definition: $\cos\theta = 1$ iff $\theta = 0$; $\cos\theta = 0$ iff $\theta = \pi/2$.
\end{proof}

\begin{example}[Two lines in $\bbR^2$]
\lean{CoherenceCalculus.AppendixOperators.lines_r2_example}
\label{ex:lines-r2}
Let $U = \Span\{e_1\}$ and $V = \Span\{\cos\theta\, e_1 + \sin\theta\, e_2\}$ for $\theta \in [0, \pi/2]$.
The single principal angle is $\theta$, and the frustration tensor is $\mathcal{F}(U,V) = (\cos\theta)$.
\end{example}

\begin{remark}[Connection to representation projectors]
\lean{CoherenceCalculus.AppendixOperators.IsotypicProperties}
If $U$ and $V$ are isotypic subspaces for different irreducible representations of a group $G$ (Section~\ref{subsec:representation-projector}), then $U \perp V$ by Schur orthogonality, so all principal angles are $\pi/2$.
For subspaces corresponding to the \emph{same} irrep, the principal angles measure how ``aligned'' different copies are.
\end{remark}

\begin{quote}\small
\textbf{Intuition.}
Principal angles quantify ``mixing'' or ``frustration'' between two constraint subspaces.
The frustration tensor is the SVD data packaging all pairwise alignments.
When subspaces come from symmetry constraints, the principal angles measure compatibility.
\end{quote}

% ------------------------------------------------------------
\subsection{Field Filters and Invariant Subalgebras (Selection by Closure)}
\label{subsec:field-filter}
% ------------------------------------------------------------

A field filter is a projection operator that enforces coefficient or structure restrictions by projecting onto an invariant subalgebra.

\begin{definition}[Algebra of observables]
\label{def:observable-algebra}\lean{CoherenceCalculus.AppendixOperators.ObservableAlgebra}
Let $\mathcal{A}$ be a real or complex $*$-algebra of bounded linear operators on $\cH_{\mathrm{obs}}$.
Elements of $\mathcal{A}$ are ``observables'' in the algebraic sense.
\end{definition}

\begin{definition}[Invariant subalgebra]
\label{def:invariant-subalgebra}\lean{CoherenceCalculus.AppendixOperators.GroupAction}\lean{CoherenceCalculus.AppendixOperators.invariantSubalgebra}
Let $\Gamma$ be a finite group acting on $\mathcal{A}$ by $*$-automorphisms: for each $g \in \Gamma$, there is a map $\alpha_g: \mathcal{A} \to \mathcal{A}$ satisfying
\begin{enumerate}[label=(\roman*)]
\item $\alpha_g(a + b) = \alpha_g(a) + \alpha_g(b)$,
\item $\alpha_g(ab) = \alpha_g(a)\alpha_g(b)$,
\item $\alpha_g(a^*) = \alpha_g(a)^*$,
\item $\alpha_{gh} = \alpha_g \circ \alpha_h$.
\end{enumerate}
The \emph{$\Gamma$-invariant subalgebra} is
\[
\mathcal{A}^\Gamma := \{ a \in \mathcal{A} : \alpha_g(a) = a \text{ for all } g \in \Gamma \}.
\]
\end{definition}

\begin{definition}[Invariant filter (group averaging)]
\label{def:invariant-filter}\lean{CoherenceCalculus.AppendixOperators.invariantFilter}
The \emph{invariant filter} associated to the group action $\Gamma$ is
\[
\mathcal{G}_\Gamma: \mathcal{A} \to \mathcal{A}, \qquad
\mathcal{G}_\Gamma(a) := \frac{1}{|\Gamma|} \sum_{g \in \Gamma} \alpha_g(a).
\]
\end{definition}

\begin{proposition}[Properties of the invariant filter]
\label{prop:invariant-filter-properties}\lean{CoherenceCalculus.AppendixOperators.invariantFilter_linear}\lean{CoherenceCalculus.AppendixOperators.invariantFilter_idempotent}\lean{CoherenceCalculus.AppendixOperators.invariantFilter_range}\lean{CoherenceCalculus.AppendixOperators.invariantFilter_fixed_points}\lean{CoherenceCalculus.AppendixOperators.invariantFilter_preserves_star}
The invariant filter $\mathcal{G}_\Gamma$ satisfies:
\begin{enumerate}[label=(\roman*)]
\item \textbf{Linearity:} $\mathcal{G}_\Gamma(\lambda a + \mu b) = \lambda \mathcal{G}_\Gamma(a) + \mu \mathcal{G}_\Gamma(b)$.
\item \textbf{Idempotence:} $\mathcal{G}_\Gamma^2 = \mathcal{G}_\Gamma$ (projection property).
\item \textbf{Range:} $\Img(\mathcal{G}_\Gamma) = \mathcal{A}^\Gamma$.
\item \textbf{Fixed points:} $\mathcal{G}_\Gamma(a) = a$ if and only if $a \in \mathcal{A}^\Gamma$.
\item \textbf{$*$-preservation:} $\mathcal{G}_\Gamma(a^*) = \mathcal{G}_\Gamma(a)^*$.
\end{enumerate}
\end{proposition}

\begin{proof}
(i) Linearity follows from linearity of each $\alpha_g$ and of summation.

(ii) For idempotence:
\[
\mathcal{G}_\Gamma^2(a) = \frac{1}{|\Gamma|} \sum_{g \in \Gamma} \alpha_g\Bigl( \frac{1}{|\Gamma|} \sum_{h \in \Gamma} \alpha_h(a) \Bigr)
= \frac{1}{|\Gamma|^2} \sum_{g,h \in \Gamma} \alpha_{gh}(a)
= \frac{1}{|\Gamma|^2} \sum_{g \in \Gamma} \sum_{k \in \Gamma} \alpha_k(a)
= \mathcal{G}_\Gamma(a),
\]
where we used the substitution $k = gh$ and the fact that as $h$ ranges over $\Gamma$, so does $k = gh$.

(iii) If $a = \mathcal{G}_\Gamma(b)$, then $\alpha_g(a) = \frac{1}{|\Gamma|} \sum_{h} \alpha_{gh}(b) = \mathcal{G}_\Gamma(b) = a$, so $a \in \mathcal{A}^\Gamma$.
Conversely, if $a \in \mathcal{A}^\Gamma$, then $\mathcal{G}_\Gamma(a) = \frac{1}{|\Gamma|} \sum_g a = a$.

(iv) Follows from (ii) and (iii).

(v) Each $\alpha_g$ preserves the $*$-operation, so $\mathcal{G}_\Gamma(a^*) = \frac{1}{|\Gamma|} \sum_g \alpha_g(a)^* = \mathcal{G}_\Gamma(a)^*$.
\end{proof}

\begin{example}[Real vs.\ complex: conjugation filter]
\lean{CoherenceCalculus.AppendixOperators.conjugation_filter_example}
\label{ex:conjugation-filter}
Let $\mathcal{A} = M_n(\mathbb{C})$ (complex $n \times n$ matrices) and $\Gamma = \{e, \sigma\}$ where $\sigma$ is complex conjugation: $\alpha_\sigma(A) = \bar{A}$.
Then
\[
\mathcal{G}_\Gamma(A) = \tfrac{1}{2}(A + \bar{A}) = \mathrm{Re}(A),
\]
and $\mathcal{A}^\Gamma = M_n(\bbR)$.
The filter projects onto real matrices.
\end{example}

\begin{example}[Diagonal filter]
\lean{CoherenceCalculus.AppendixOperators.diagonal_filter_example}
\label{ex:diagonal-filter}
Let $\mathcal{A} = M_n(\mathbb{C})$ and let $\Gamma = (\mathbb{Z}/n\mathbb{Z})$ act by cyclic permutation of basis vectors.
Then $\mathcal{A}^\Gamma$ consists of circulant matrices.
For the group of all permutations $\Gamma = S_n$, the invariant subalgebra is scalar multiples of the identity.
\end{example}

\begin{remark}[Connection to spectral selector]
\lean{CoherenceCalculus.AppendixOperators.spectralSelector}
The invariant filter $\mathcal{G}_\Gamma$ and the spectral selector $Q_S$ (Section~\ref{subsec:quantization}) are both selection operators.
$Q_S$ selects eigenspaces; $\mathcal{G}_\Gamma$ selects by symmetry invariance.
Both are idempotent and extract a canonical ``admissible'' component.
\end{remark}

\begin{quote}\small
\textbf{Intuition.}
A field/invariant filter enforces coefficient or structure restrictions by averaging over a symmetry group.
The result is a projection onto the subalgebra of invariants, providing a
direct way to extract the admissible part of any observable.
\end{quote}

% ------------------------------------------------------------
\subsection{Unitary Splittings and Unity Identities}
\label{subsec:unity-splitting}
% ------------------------------------------------------------

Unitary splittings provide norm bookkeeping: total norm partitions across observable, shadow, and defect components.

\begin{definition}[Observable-shadow decomposition]
\label{def:obs-shadow-decomp}\lean{CoherenceCalculus.AppendixOperators.obsShadowDecomp}
For a Hilbert space $\cH$ with horizon projection $\CCPi{h}$ and shadow projection $\CCPibar{h} = I - \CCPi{h}$, every vector $x \in \cH$ decomposes as
\[
x = \CCPi{h} x + \CCPibar{h} x = x_{\mathrm{obs}} + x_{\mathrm{sh}},
\]
where $x_{\mathrm{obs}} := \CCPi{h} x \in \cH_{\mathrm{obs}}$ and $x_{\mathrm{sh}} := \CCPibar{h} x \in \cH_{\mathrm{sh}}$.
\end{definition}

\begin{proposition}[Pythagorean splitting of norms]
\label{prop:pythagorean-norm}\lean{CoherenceCalculus.AppendixOperators.pythagorean_splitting}
For the orthogonal decomposition $x = x_{\mathrm{obs}} + x_{\mathrm{sh}}$:
\[
\|x\|^2 = \|x_{\mathrm{obs}}\|^2 + \|x_{\mathrm{sh}}\|^2 = \|\CCPi{h} x\|^2 + \|\CCPibar{h} x\|^2.
\]
\end{proposition}

\begin{proof}
Since $\CCPi{h}$ and $\CCPibar{h}$ are orthogonal projections with $\CCPi{h} \CCPibar{h} = 0$, we have $\langle x_{\mathrm{obs}}, x_{\mathrm{sh}} \rangle = 0$.
Thus $\|x\|^2 = \langle x, x \rangle = \|x_{\mathrm{obs}}\|^2 + 2\langle x_{\mathrm{obs}}, x_{\mathrm{sh}} \rangle + \|x_{\mathrm{sh}}\|^2 = \|x_{\mathrm{obs}}\|^2 + \|x_{\mathrm{sh}}\|^2$.
\end{proof}

\begin{definition}[Defect energy]
\label{def:defect-energy-unity}\lean{CoherenceCalculus.AppendixOperators.defectEnergy_eq_leakage_norm}
For a nilpotent differential $d$ with leakage operator $\CCLeak{d}{h} = \CCLeakExpr{d}{h}$, the \emph{defect energy} of $x \in \cH_{\mathrm{obs}}$ is
\[
\CCDefEnergy{h}{x} := \|\CCLeak{d}{h} x\|^2 = \langle x, \CCLeak{d}{h}^\ast \CCLeak{d}{h}\, x \rangle.
\]
\end{definition}

\begin{proposition}[Defect energy decomposition]
\label{prop:defect-decomp}\lean{CoherenceCalculus.AppendixOperators.defect_energy_decomposition}
For $x \in \cH_{\mathrm{obs}}$ and $y := dx$ (the result of applying $d$), write $y = y_{\mathrm{obs}} + y_{\mathrm{sh}}$ where $y_{\mathrm{obs}} = \CCPi{h} y = \CCObs{d}{h} x$ and $y_{\mathrm{sh}} = \CCPibar{h} y = \CCLeak{d}{h} x$.
Then:
\[
\|dx\|^2 = \|\CCObs{d}{h} x\|^2 + \CCDefEnergy{h}{x}.
\]
That is, the squared norm of $dx$ splits into the observed part $\|\CCObs{d}{h} x\|^2$ and the defect $\CCDefEnergy{h}{x}$.
\end{proposition}

\begin{proof}
Apply Proposition~\ref{prop:pythagorean-norm} to $y = dx$:
\[
\|dx\|^2 = \|\CCPi{h} dx\|^2 + \|\CCPibar{h} dx\|^2 = \|\CCObs{d}{h} x\|^2 + \|\CCLeak{d}{h} x\|^2 = \|\CCObs{d}{h} x\|^2 + \CCDefEnergy{h}{x}. \qedhere
\]
\end{proof}

\begin{corollary}[Norm bound on observed operator]
\label{cor:norm-bound}\lean{CoherenceCalculus.AppendixOperators.norm_bound_observed_statement}\lean{CoherenceCalculus.AppendixOperators.norm_bound_equality_iff}
For $x \in \cH_{\mathrm{obs}}$:
\[
\|\CCObs{d}{h} x\|^2 \le \|dx\|^2, \qquad \text{with equality iff } \CCDefEnergy{h}{x} = 0.
\]
The observed operator cannot increase squared norm beyond the full operator; defect accounts for any ``missing'' norm.
\end{corollary}

\begin{proposition}[Isometry property from polar decomposition]
\label{prop:isometry-property}\lean{CoherenceCalculus.AppendixOperators.PolarDecomp}\lean{CoherenceCalculus.AppendixOperators.isometry_property}\lean{CoherenceCalculus.AppendixOperators.leakage_norm_eq_abs_norm}
Let $\CCLeak{d}{h} = U_h |\CCLeak{d}{h}|$ be the polar decomposition of the leakage operator (Section~\ref{subsec:aliasing}).
The partial isometry $U_h$ satisfies:
\[
\|U_h z\| = \|z\| \quad \text{for all } z \in \overline{\Img(|\CCLeak{d}{h}|)}.
\]
In particular, $\|\CCLeak{d}{h} x\| = \||\CCLeak{d}{h}| x\|$ for all $x$.
\end{proposition}

\begin{proof}
By definition of partial isometry, $U_h^* U_h$ is the projection onto the initial space $\overline{\Img(|\CCLeak{d}{h}|)}$.
For $z$ in this subspace, $\|U_h z\|^2 = \langle U_h z, U_h z \rangle = \langle z, U_h^* U_h z \rangle = \langle z, z \rangle = \|z\|^2$.
The second statement follows since $\CCLeak{d}{h} x = U_h |\CCLeak{d}{h}| x$ and $|\CCLeak{d}{h}| x \in \Img(|\CCLeak{d}{h}|)$.
\end{proof}

\begin{remark}[Connection to telescoping identity]
\lean{CoherenceCalculus.AppendixOperators.pythagorean_splitting}
The norm splitting $\|x\|^2 = \|\CCPi{h} x\|^2 + \|\CCPibar{h} x\|^2$ underlies the telescoping decomposition in the Horizon Fundamental Theorem (Section~\ref{sec:hft}).
Each horizon level $h$ contributes its share of the total norm, with no overlap due to orthogonality.
\end{remark}

\begin{example}[Norm accounting in $\bbR^3$]
\lean{CoherenceCalculus.AppendixOperators.norm_r3_example}
\label{ex:norm-r3}
Let $\cH = \bbR^3$ with $\CCPi{h} = e_1 e_1^* + e_2 e_2^*$ (projection onto the $(x,y)$-plane).
For $x = (a, b, c)^\top$:
\[
\|x\|^2 = a^2 + b^2 + c^2 = \underbrace{(a^2 + b^2)}_{\|\CCPi{h} x\|^2} + \underbrace{c^2}_{\|\CCPibar{h} x\|^2}.
\]
The ``shadow'' component $c^2$ is exactly what the planar observer cannot account for.
\end{example}

\begin{quote}\small
\textbf{Intuition.}
Unity identities are bookkeeping laws: total norm (or squared norm) splits cleanly across observable and shadow components.
Defect energy $\CCDefEnergy{h}{x}$ measures the portion of $\|dx\|^2$ that ``leaks'' out of the observable.
Isometries (from polar decomposition) preserve norms, so they transfer magnitude without creating or destroying it.
\end{quote}
